\documentclass[11pt,a4paper]{article}

\usepackage[T1]{fontenc}
\usepackage[utf8]{inputenc}
\usepackage{amsmath}
\usepackage{amssymb}
\usepackage{booktabs}
\usepackage{geometry}
\usepackage{hyperref}
\usepackage{microtype}

\geometry{margin=24mm}
\hypersetup{
  colorlinks=true,
  linkcolor=black,
  urlcolor=black
}

\newcommand{\R}{\mathbb{R}}
\newcommand{\clip}{\operatorname{clip}}
\newcommand{\sgn}{\operatorname{sgn}}
\newcommand{\Normal}{\mathcal{N}}

\title{AGH Racing DV Simulator\\Model Equations}
\author{AGH Racing}
\date{}

\begin{document}
\maketitle
\tableofcontents

\section{Scope and notation}

This document is the mathematical reference for the models implemented by the
simulator. The C++ implementation remains the executable definition. All
quantities use SI units. Wheel indices are
\[
  i\in\{\mathrm{FL},\mathrm{FR},\mathrm{RL},\mathrm{RR}\}.
\]
The vehicle body frame uses longitudinal velocity \(v_x\), lateral velocity
\(v_y\), yaw angle \(\psi\), and yaw rate \(r=\dot{\psi}\).

\section{Planar four-wheel vehicle dynamics}

Tyre forces are calculated in each wheel frame and rotated into the body frame.
The body equations are
\begin{align}
  \dot v_x &= \frac{\sum_i F_{x,i}^{b}-F_{\mathrm{res}}}{m}+v_y r,\\
  \dot v_y &= \frac{\sum_i F_{y,i}^{b}}{m}-v_x r,\\
  \dot r   &= \frac{M_z}{I_z}.
\end{align}
The yaw moment follows directly from the wheel positions
\begin{equation}
  M_z=\sum_i\left(x_iF_{y,i}^{b}-y_iF_{x,i}^{b}\right).
\end{equation}
Global pose kinematics are
\begin{align}
  \dot x &= v_x\cos\psi-v_y\sin\psi,\\
  \dot y &= v_x\sin\psi+v_y\cos\psi,\\
  \dot\psi &= r.
\end{align}

The longitudinal resistance model is
\begin{align}
  F_{\mathrm{roll}} &= C_rmg,\\
  F_{\mathrm{drag}} &= C_d|v_x|^2,\\
  F_{\mathrm{extra}} &= F_0+k_v|v_x|,\\
  F_{\mathrm{res}} &=
  \sgn(v_x)\left(F_{\mathrm{roll}}+F_{\mathrm{drag}}+
  F_{\mathrm{extra}}\right).
\end{align}
Aerodynamic downforce is split between axles:
\begin{equation}
  F_{\mathrm{aero},F}=C_{l,F}v_x^2,\qquad
  F_{\mathrm{aero},R}=C_{l,R}v_x^2.
\end{equation}

Each wheel has an independent angular-velocity state
\begin{equation}
  I_w\dot\omega_i=T_i-RF_{x,i}.
\end{equation}

\section{Control modes and torque routing}

The two values of \texttt{dv\_interfaces/Control.move\_type} use separate
paths.

\subsection{\texttt{ONE\_WHEEL}}

Only this mode uses baseline allocation and torque vectoring. The total motor
torque request is
\begin{equation}
  T_{\Sigma}^{m} =
  T_{\mathrm{FL}}^{m}+T_{\mathrm{FR}}^{m}+
  T_{\mathrm{RL}}^{m}+T_{\mathrm{RR}}^{m}.
\end{equation}
After the fixed drivetrain ratio \(i_g\), baseline wheel torques are
\begin{align}
  T_{\mathrm{FL}}^{0}=T_{\mathrm{FR}}^{0}
    &=\frac{1}{2}\lambda_F i_gT_{\Sigma}^{m},\\
  T_{\mathrm{RL}}^{0}=T_{\mathrm{RR}}^{0}
    &=\frac{1}{2}(1-\lambda_F)i_gT_{\Sigma}^{m},
\end{align}
where \(\lambda_F\) has separate configured values for drive and braking.
Torque vectoring applies a bounded symmetric correction
\begin{equation}
  \begin{bmatrix}
  T_{\mathrm{FL}}\\T_{\mathrm{FR}}\\T_{\mathrm{RL}}\\T_{\mathrm{RR}}
  \end{bmatrix}
  =
  \begin{bmatrix}
  T_{\mathrm{FL}}^{0}\\T_{\mathrm{FR}}^{0}\\
  T_{\mathrm{RL}}^{0}\\T_{\mathrm{RR}}^{0}
  \end{bmatrix}
  +
  \begin{bmatrix}
  +\Delta T\\-\Delta T\\+\Delta T\\-\Delta T
  \end{bmatrix}.
\end{equation}

\subsection{\texttt{FOUR\_WHEEL}}

Baseline allocation and torque vectoring are bypassed. Each request is
converted independently:
\begin{equation}
  T_i=i_gT_i^m.
\end{equation}

\section{MF6.1 tyre model}

\subsection{Slip variables}

For local wheel-frame velocities \(v_{x,i}\), \(v_{y,i}\),
\begin{align}
  v_{\mathrm{den}}(v_x) &=
  \begin{cases}
  |v_x|, & |v_x|>v_{\mathrm{th}},\\
  \dfrac{1}{2}\left(v_{\mathrm{th}}+
  \dfrac{v_x^2}{v_{\mathrm{th}}}\right), & |v_x|\leq v_{\mathrm{th}},
  \end{cases}\\
  \kappa_i &=
  \clip\left(\frac{\omega_iR-v_{x,i}}{v_{\mathrm{den}}(v_{x,i})},
  -0.99,0.99\right),\\
  \alpha_i &=-\operatorname{atan2}(v_{y,i},v_{x,i,\mathrm{safe}}).
\end{align}
The implementation uses \(v_{\mathrm{th}}=1.5\,\mathrm{m\,s^{-1}}\).

\subsection{Pure longitudinal and lateral force}

For either longitudinal slip or slip angle, the common Magic Formula form is
\begin{align}
  \phi &= Bx-E\left(Bx-\arctan(Bx)\right),\\
  F(x) &= D\sin\left(C\arctan\phi\right).
\end{align}
Load sensitivity uses
\begin{equation}
  df_z=\frac{F_z-F_{z0}}{F_{z0}},
\end{equation}
with peak factors
\begin{equation}
  \mu_x=\lambda_x(p_{Dx1}+p_{Dx2}df_z),\qquad
  \mu_y=\lambda_y(p_{Dy1}+p_{Dy2}df_z),
\end{equation}
and
\begin{equation}
  D_x=\mu_xF_z,\qquad D_y=\mu_yF_z.
\end{equation}
The longitudinal stiffness term is
\begin{equation}
  K_{x\kappa}=
  F_z(p_{Kx1}+p_{Kx2}df_z)\exp(p_{Kx3}df_z),
  \qquad B_x=\frac{K_{x\kappa}}{C_xD_x}.
\end{equation}
The lateral stiffness term is
\begin{equation}
  K_{y\alpha}=p_{Ky1}F_{z0}
  \sin\left(p_{Ky4}\arctan
  \frac{F_z}{p_{Ky2}F_{z0}}\right),
  \qquad B_y=\frac{K_{y\alpha}}{C_yD_y}.
\end{equation}
Horizontal shifts are applied before evaluating the curves, and the force
value at zero physical slip is subtracted to remove the shift-induced offset.

\subsection{Combined slip and force relaxation}

Pure forces are limited by a friction ellipse. Define
\begin{equation}
  s=\sqrt{
  \left(\frac{F_{x0}}{\mu_xF_z}\right)^2+
  \left(\frac{F_{y0}}{\mu_yF_z}\right)^2}.
\end{equation}
Then
\begin{equation}
  (F_x,F_y)=
  \begin{cases}
  (F_{x0},F_{y0}), & s\leq1,\\
  (F_{x0}/s,F_{y0}/s), & s>1.
  \end{cases}
\end{equation}

With \(V=\sqrt{v_x^2+v_{\mathrm{th}}^2}\), the dynamic relaxation lengths are
\begin{align}
  L_\alpha &= \frac{Vd_y}{c_y}
  +\frac{1}{c_y}\left|\frac{\partial F_y}{\partial\alpha}\right|,\\
  L_\kappa &= \frac{Vd_x}{c_x}
  +\frac{1}{c_x}\left|\frac{\partial F_x}{\partial\kappa}\right|.
\end{align}
For a frozen target over one integration substep,
\begin{equation}
  F^{+}=F+(F_{\mathrm{MF}}-F)
  \left(1-\exp\left[-\frac{V\Delta t}{L}\right]\right).
\end{equation}

\section{Suspension and normal loads}

Each wheel normal load is decomposed as
\begin{equation}
  N_i=N_{\mathrm{direct},i}
  +N_{\mathrm{lat,el},i}
  +N_{\mathrm{long,el},i}.
\end{equation}
The direct component contains static load, aerodynamic downforce, unsprung
transfer, and geometric transfer.

For \(L=l_f+l_r\), static axle loads are
\begin{equation}
  F_{z,F}^{0}=\frac{mgl_r}{L},\qquad
  F_{z,R}^{0}=\frac{mgl_f}{L}.
\end{equation}
Longitudinal unsprung and geometric transfer are
\begin{align}
  \Delta F_{z,\mathrm{long},u}
  &=\frac{a_x(m_{uF}h_{uF}+m_{uR}h_{uR})}{L},\\
  \Delta F_{z,\mathrm{long},g}
  &=k_{gF}F_{x,F,s}+k_{gR}F_{x,R,s}.
\end{align}
The elastic longitudinal target is
\begin{equation}
  \Delta F_{z,\mathrm{long},el}^{*}
  =\frac{m_sa_xh_s}{L}
  -\Delta F_{z,\mathrm{long},g}.
\end{equation}

Per-axle lateral direct transfer is
\begin{align}
  \Delta F_{z,u,j} &= \frac{m_{u,j}a_yh_{u,j}}{t_j},\\
  \Delta F_{z,g,j} &= \frac{F_{y,j,s}h_{\mathrm{RC},j}}{t_j},
  \qquad j\in\{F,R\}.
\end{align}
The remaining sprung roll moment is
\begin{equation}
  M_{\mathrm{roll},el}^{*}
  =m_sa_yh_s-
  \left(F_{y,F,s}h_{\mathrm{RC},F}
  +F_{y,R,s}h_{\mathrm{RC},R}\right).
\end{equation}
It is split between axles:
\begin{equation}
  \Delta F_{z,el,F}^{*}
  =\frac{\lambda_\phi M_{\mathrm{roll},el}^{*}}{t_F},
  \qquad
  \Delta F_{z,el,R}^{*}
  =\frac{(1-\lambda_\phi)M_{\mathrm{roll},el}^{*}}{t_R}.
\end{equation}
Elastic pitch and roll loads use independent second-order dynamics:
\begin{equation}
  \ddot N_{\mathrm{el}}
  =\omega^2(N_{\mathrm{el}}^{*}-N_{\mathrm{el}})
  -2\zeta\omega\dot N_{\mathrm{el}}.
\end{equation}
The previous-step acceleration is used to avoid the algebraic loop
\(N\rightarrow F_{\mathrm{tyre}}\rightarrow a\rightarrow N\).

\section{Steering and drivetrain}

The steering rack is a second-order system:
\begin{equation}
  \ddot\delta_r=
  \omega_n^2(\delta_{\mathrm{request}}-\delta_r)
  -2\zeta\omega_n\dot\delta_r.
\end{equation}
Left and right wheel angles are obtained from the rack geometry. The command
and measured encoder include independent configured noise terms; the actuator
command additionally includes steering bias.

Each drivetrain torque state follows
\begin{equation}
  \dot T_i=\frac{T_{i,\mathrm{request}}-T_i}{\tau_d}.
\end{equation}
Per-wheel torque is limited by the configured torque bounds and by one quarter
of the total drive or regenerative power limit.

\section{Camera model}

For a cone at global coordinates \((x_c,y_c,z_c)\), the planar transformation
to the vehicle frame is
\begin{equation}
  \begin{bmatrix}x_b\\y_b\end{bmatrix}
  =
  \begin{bmatrix}
  \cos\psi&\sin\psi\\
  -\sin\psi&\cos\psi
  \end{bmatrix}
  \begin{bmatrix}x_c-x\\y_c-y\end{bmatrix}.
\end{equation}
Sensor offsets are added to obtain \((x_s,y_s,z_s)\). A cone is visible when
\begin{equation}
  x_s>0,\qquad d\leq d_{\max},\qquad
  |y_s|\leq x_s\tan\frac{FOV_h}{2},\qquad
  |z_s|\leq x_s\tan\frac{FOV_v}{2}.
\end{equation}
Position noise grows exponentially with distance:
\begin{align}
  RMSE(d)&=a\exp(bd),\\
  \sigma_{xy}&=\frac{RMSE(d)}{\sqrt2},\\
  x_{\mathrm{obs}}&=x_s+\sigma_{xy}\eta_x,\qquad
  y_{\mathrm{obs}}=y_s+\sigma_{xy}\eta_y,
\end{align}
where \(\eta_x,\eta_y\sim\Normal(0,1)\).

\section{Lidar model}

After the body transform and sensor translation, lidar pitch is applied:
\begin{align}
  x_l &= \cos\theta_l\,x_{\mathrm{rel}}
       +\sin\theta_l\,z_{\mathrm{rel}},\\
  y_l &= y_{\mathrm{rel}},\\
  z_l &= -\sin\theta_l\,x_{\mathrm{rel}}
       +\cos\theta_l\,z_{\mathrm{rel}}.
\end{align}
The ideal planar measurement is
\begin{equation}
  r=\sqrt{x_l^2+y_l^2},\qquad
  \varphi=\operatorname{atan2}(y_l,x_l).
\end{equation}
The visibility gate is
\begin{equation}
  x_l>0,\qquad r\leq r_{\max},\qquad
  |\varphi|\leq\frac{\Phi_{\mathrm{ROI}}}{2}.
\end{equation}
Optional non-deskewed motion distortion contributes
\begin{equation}
  \sigma_{\varphi,\mathrm{motion}}
  =\frac{|r_{\mathrm{yaw}}|T_{\mathrm{ROI}}}{\sqrt{12}},
\end{equation}
and total azimuth uncertainty is
\begin{equation}
  \sigma_{\varphi}=
  \sqrt{\sigma_{\varphi,0}^2+
  \sigma_{\varphi,\mathrm{motion}}^2}.
\end{equation}
The observed polar coordinates are
\begin{equation}
  r_{\mathrm{obs}}=r+\sigma_r\eta_r,\qquad
  \varphi_{\mathrm{obs}}=\varphi+\sigma_\varphi\eta_\varphi,
  \qquad \eta_r,\eta_\varphi\sim\Normal(0,1).
\end{equation}
They are converted back to Cartesian coordinates:
\begin{equation}
  x_{\mathrm{obs}}=r_{\mathrm{obs}}\cos\varphi_{\mathrm{obs}},
  \qquad
  y_{\mathrm{obs}}=r_{\mathrm{obs}}\sin\varphi_{\mathrm{obs}}.
\end{equation}

\section{Fusion and perception errors}

Fusion uses lidar position. Camera colour is attached only when the same cone
passes the camera visibility gate. Lidar-only cones have unknown colour.

Optional dropout applies an independent Bernoulli trial:
\begin{equation}
  D_i\sim\operatorname{Bernoulli}(p_{\mathrm{drop}}).
\end{equation}
A detection is removed when \(D_i=1\). The number of optional false positives
per frame is
\begin{equation}
  N_{\mathrm{false}}\sim\operatorname{Poisson}(\lambda_{\mathrm{false}}).
\end{equation}
Their forward range and lateral position are sampled uniformly inside the
configured region.

Camera and lidar execution latency are sampled as
\begin{equation}
  t_{\mathrm{exec}}=
  \left|\Normal(\mu_{\mathrm{exec}},\sigma_{\mathrm{exec}}^2)\right|.
\end{equation}

\section{IMU}

Accelerometer and gyroscope biases are independent random walks:
\begin{equation}
  b_{k+1}=b_k+\sigma_{\mathrm{rw}}\sqrt{\Delta t}\,\eta_k,
  \qquad \eta_k\sim\Normal(0,1).
\end{equation}
Each channel uses a first-order bandwidth model:
\begin{equation}
  \tau=\frac{1}{2\pi f_c},\qquad
  \alpha=\frac{\Delta t}{\tau+\Delta t},\qquad
  y_k^f=y_{k-1}^f+\alpha(y_k-y_{k-1}^f).
\end{equation}
The sampled output is
\begin{equation}
  y_k^{m}=y_k^f+b_k+\sigma_n\eta_k.
\end{equation}
Accelerometer and gyroscope have independent acquisition clocks. Output
messages contain the arithmetic mean of the configured number of buffered
samples.

\section{State estimator}

For each estimated state component \(q\),
\begin{align}
  b_{q,k+1}&=b_{q,k}+\sigma_{q,\mathrm{rw}}
  \sqrt{\Delta t}\,\eta_{b,k},\\
  \hat q_k&=q_k+b_{q,k}+\sigma_{q,n}\eta_{n,k}.
\end{align}
Position and yaw use one delayed queue; body velocities use another. The
estimate is published only after the configured calibration interval and when
both delayed samples are available.

\section{Discrete timing}

For device period \(T_d\) and physics step \(\Delta t\),
\begin{equation}
  N_d=\max\left(1,\operatorname{round}\frac{T_d}{\Delta t}\right).
\end{equation}
A device runs when
\begin{equation}
  k\bmod N_d=k_{\mathrm{phase}}.
\end{equation}
Independent randomized phases avoid artificial simultaneous sampling.
Sampled values are held in caches and computation or transport latency is
represented by queues indexed by their ready simulation step.

\section{Implementation references}

\begin{center}
\begin{tabular}{ll}
\toprule
Model & Source\\
\midrule
Body and wheel dynamics & \texttt{src\_helpers/double\_track.cpp}\\
Tyres & \texttt{src\_helpers/tire\_model.cpp}\\
Suspension & \texttt{src\_helpers/suspension.cpp}\\
Steering & \texttt{src\_helpers/steering\_system.cpp}\\
Torque routing & \texttt{src\_helpers/torque\_allocation.cpp}\\
Camera, lidar, fusion & \texttt{src\_helpers/cone\_detector.cpp}\\
IMU & \texttt{src\_helpers/imu\_simulator.cpp}\\
Estimator and timing & \texttt{sim\_loop.cpp}\\
\bottomrule
\end{tabular}
\end{center}

\end{document}
